% Abstract C of ABSTRACTS.md, trimmed a fifth, with the headline t folded in.
\begin{abstract}
\noindent Trademark goods/services descriptions are a forced, dated
disclosure of what a firm expected to sell, written by the millions of
firms that never patent. This paper builds two measures on them and tests
what each predicts. A description is reduced to themes by a topic model
fitted once over the corpus; its theme mix is compared, by
Kullback--Leibler divergence, to its own Nice class in the five years
before and after filing. The average of the two divergences is
atypicality; their difference is lead. Three partitions of the vocabulary
(fifty global themes, five hundred, fifty per class), two reference
weightings, and five window constructions are compared on identical
filings. Lead is robust to all of them and predicts cancellation for
non-use at the five-year proof of continued use: $+2.3$, $+2.6$ and
$+1.9$ points under the three partitions, $+1.4$ with design controls
($t = 8.3$). Atypicality is not: its protective effect at the proof is
$4.4$ points under global themes and $0.7$ under per-class themes, and
its sign at the listing margin depends on how description length is held.
Two hazards are documented for trademark-text research. Scoring words
rather than theme distributions measures the length of the filing; and
the registered description is a conformed document, because applicants
who refile rewrite toward their class's typical language. The measure's
validity is therefore strongest for lead, for outcomes at the mark, and
within an industry; the paper states where each condition fails.
\end{abstract}
